\documentclass[12pt,a4paper]{article}

\usepackage{fontspec}
\usepackage[margin=2cm]{geometry}
\usepackage{enumitem}
\usepackage{booktabs}
\usepackage{tabularx}
\usepackage[table]{xcolor}
\usepackage{titlesec}
\usepackage[hidelinks]{hyperref}

\setlist{nosep, leftmargin=*}
\titlespacing*{\section}{0pt}{1.0ex}{0.5ex}
\titlespacing*{\subsection}{0pt}{0.8ex}{0.3ex}
\titleformat{\section}{\large\bfseries}{\thesection}{0.6em}{}
\titleformat{\subsection}{\normalsize\bfseries}{}{0em}{}
\setlength{\parindent}{0pt}
\setlength{\parskip}{0.4ex}
\pagestyle{plain}

% Visible placeholder for anything the group still has to fill in.
\newcommand{\todo}[1]{\textcolor{red}{[TBD: #1]}}
\newcommand{\appname}{SUnite}

\begin{document}

\begin{center}
  {\large\bfseries Sabancı University}\\[0.6ex]
  {\large\bfseries CS310 -- Mobile Application Development}\\[0.6ex]
  {\large\bfseries 2026-2027 Fall Project Step 1 Proposal}\\[2.5ex]
  {\LARGE\bfseries \appname{}} 
  Find your group at Sabancı\\[2ex]
  {\bfseries Group 32}\\[1.5ex]
  Azra At \quad 34090\\
  Ece Ulutaş \quad 35589\\
  Elif Çamlıkaya \quad 35191\\
  Mert Ali İlter \quad 34447\\[1.5ex]
  GitHub: \url{https://github.com/ssinix/sunite}
\end{center}

\section{App Overview}

\textbf{What the app does:} \appname{} is a mobile app, exclusive to Sabancı
University students, for finding and forming groups of any kind. It is
organised as a set of category boards on which students can post and browse
listings of two kinds: \emph{existing groups looking for new members}, and
\emph{students looking for a group to join}. Boards cover course groups,
Erasmus, startups/projects and campus events.

\textbf{Problem it addresses:} Finding the right people at SU is currently
scattered across WhatsApp groups, SUCourse forums and word of
mouth. A student who needs a project partner for a course, wants to meet others
going to the same Erasmus school, is looking for co-founders with a specific
skill, or simply does not want to go to a club event alone has no single place
to look. \appname{} gathers all of these requests in one place and makes them
searchable by course, school, skill and event.

\textbf{Intended users:} Sabancı University undergraduate and graduate
students, verified through their \texttt{@sabanciuniv.edu} e-mail address.
Clubs and student initiatives benefit indirectly, since students can organise
groups around their events.

\section{Core Features}

\begin{enumerate}
  \item \textbf{SU-only authentication:} sign up and log in with a
        university e-mail address, with e-mail verification.
  \item \textbf{User profiles:} name, e-mail, year of study, major(s),
        minor(s), skills/specialties and current groups; users can view each
        other's profiles.
  \item \textbf{Four category boards:} \emph{Course Groups},
        \emph{Erasmus}, \emph{Startups \& Projects} and \emph{Events}.
  \item \textbf{Two post types:} ``group looking for members'' and
        ``person looking for a group''.
  \item \textbf{Create, edit and delete posts/groups:} title, description,
        capacity (current/maximum), ``looking for'' (specific qualities or
        skills needed) and a category-specific field (course code, Erasmus
        school and term, or event name and date).
  \item \textbf{Browse, search and filter:} filter a board by course code,
        school, tag or groups that still have open spots.
  \item \textbf{Join requests:} a student requests to join a group; the group
        owner accepts or declines. Capacity updates automatically and the post
        closes when the group is full.
  \item \textbf{My Groups page:} groups the user owns or belongs to, their
        members and pending join requests.
\end{enumerate}

\subsection{Optional Features (if time permits)}
\begin{itemize}
  \item \textbf{Random matchmaking:} ``match me'' with other students who
        are also looking for a group for the same course or event.
  \item \textbf{Friends:} add friends and see which groups they are in.
  \item \textbf{In-app group chat} for members of a group.
  \item \textbf{Push notifications} for new join requests and acceptances.
  \item \textbf{Course list autocomplete} from the SU course catalogue.
\end{itemize}

\section{Data Requirements}

The app requires a cloud database so that posts and groups are shared between
all users in real time. We plan to use \textbf{Firebase Authentication} for
accounts and \textbf{Cloud Firestore} (NoSQL) for application data. No external
API is required for the core features. Small local preferences (e.g.\ last used
filters) will be kept on the device with \texttt{shared\_preferences}.

\begin{tabularx}{\textwidth}{@{}lX@{}}
  \toprule
  \textbf{Collection} & \textbf{Main fields} \\
  \midrule
  \texttt{users}        & uid, name, e-mail, year, majors, minors, skills, group IDs \\
  \texttt{groups}       & title, description, category, capacity, owner ID, member IDs,
                          ``looking for'', category field (course code / school \& term / event \& date) \\
  \texttt{posts}        & type (group seeking members / person seeking group), category,
                          author ID, linked group ID, text, created at, status (open/closed) \\
  \texttt{joinRequests} & group ID, user ID, status (pending/accepted/declined), created at \\
  \midrule
  \multicolumn{2}{@{}l}{\emph{Optional:} \texttt{friendships} (user pairs), \texttt{messages} (group chat).} \\
  \bottomrule
\end{tabularx}

\section{Potential Challenges}

\begin{itemize}
  \item \textbf{Restricting access to SU students:} enforcing the e-mail
        domain and verification, and writing Firestore security rules so users
        can only edit their own data.
  \item \textbf{Data modelling and queries:} Firestore has limits on
        multi-field filtering, so board filters will need composite indexes
        and some data duplication.
  \item \textbf{Consistency:} two students may be accepted into the last spot
        at the same time; membership and capacity updates must use
        transactions.
  \item \textbf{State management:} keeping boards, group pages and profiles
        in sync with real-time updates across many screens.
  \item \textbf{Privacy and moderation:} letting users control which profile
        details are visible and handling inappropriate posts.
\end{itemize}

\section{Team Responsibilities}

\textbf{Step 1 coordinator:} Mert Ali İlter 

% Suggested task buckets for later assignment:
%   - Writing proposal sections (overview, features, data, challenges)
%   - Creating the GitHub repo, README.md and Flutter .gitignore; initial commit
%   - Adding all members, the instructor and the TA as collaborators
%   - Proofreading, compiling and submitting the PDF on SUCourse
%   - Organising meetings and communicating with the TA (coordinator)
\vspace{0.5ex}
\begin{tabularx}{\textwidth}{@{}lX@{}}
  \toprule
  \textbf{Member} & \textbf{Assigned tasks} \\
  \midrule
  Azra At & Login Page and Database \\
  Ece Ulutas & Groups Page and Design \\
  Elif Çamlıkaya & Main Page and Design \\
  Mert Ali İlter & Backend \\

  \bottomrule
\end{tabularx}

\end{document}
